% Chapter 1: 执行摘要
\section{一句话结论}

英伟达 6/4 GTC 发布的 Cosmos 3 是 AI 从"生成内容"跨入"理解物理世界"的里程碑，对应制造与物流产业重塑空间约 \textbf{50 万亿美元}\footnote{NVIDIA GTC 2026 Taipei 黄仁勋主题演讲白皮书[A3]}；A股投资可按 \textbf{L0--L5 六层产业链 + 机器人/智驾双路径} 框架分层下注，首选「机构友好 + 估值有空间 + 基本面绑定深」三重共振标的，规避 7--9 月解禁高峰的游资票与 Tesla/NVIDIA 叙事错配对冲票。

\textit{数据基准：2026Q1一季报 + GTC 2026（2026.6.4）公开信息 + 6/9-6/10最新股价。}

\section{核心框架}

Cosmos 3 本身不直接变现，通过两条路径落地：
\begin{enumerate}[leftmargin=2em]
  \item \textbf{具身智能}：Cosmos 3 $\rightarrow$ Isaac Sim $\rightarrow$ GR00T $\rightarrow$ Jetson Thor $\rightarrow$ 人形机器人整机
  \item \textbf{高阶智驾}：Cosmos 3 $\rightarrow$ Omniverse/NuRec $\rightarrow$ DRIVE 软件栈 $\rightarrow$ DRIVE Thor $\rightarrow$ L3+量产车
\end{enumerate}

两条路径共享 L0 生态枢纽（协创/中科创达）与 L1 算力底座（工业富联/光模块），在 L2--L5 分道扬镳。

\section{5 星核心标的（6 家，评级上限）}

\begin{table}[H]
\centering
\small
\renewcommand{\arraystretch}{1.3}
\begin{tabularx}{\textwidth}{l l l X c}
\toprule
\textbf{标的} & \textbf{代码} & \textbf{卡位} & \textbf{确定性证据} & \textbf{Q1业绩} \\
\midrule
\textbf{协创数据} & 300857 & L0 生态枢纽 & Jetson+DRIVE+NCP 三重授权 + Cosmos Coalition 国内独家 & 净利 +343\% \\
\textbf{五一世界} & 06651.HK & L2 智驾仿真 & Omniverse生态首批合作伙伴 & 港股 \\
\textbf{索辰科技} & 688507 & L2 物理仿真 & "天工·开物"接入 Cosmos 3 推理框架 & 亏损 \\
\textbf{奥比中光} & 688322 & L3 感知入口 & Jetson 官方 3D 视觉 + Cosmos 3 适配 & 扣非 +531\% \\
\textbf{天准科技} & 688003 & L4 机器人域控 & Jetson 精英伙伴 + Thor 控制器量产 & 营收 +76\% / 扭亏 \\
\textbf{德赛西威} & 002920 & L4 智驾域控 & DRIVE Thor 首发权 + 域控市占 33.6\% & 净利 -21\% \\
\bottomrule
\end{tabularx}
\caption{5星核心标的一览}
\end{table}

\section{估值结论}

经 PEG / 季节性 / 传导系数 / 机构约束 四维校准后最终综合空间：

\begin{itemize}[leftmargin=2em]
  \item \gmark{} \textbf{首选 / 最大预期差}：德赛西威（+107\textasciitilde142\%，A 级机构友好 + Thor 量产 240 亿营收弹性）
  \item \gmark{} \textbf{上调 / 低吸}：协创数据（+21\textasciitilde49\%，PEG 0.2）、工业富联（+44\textasciitilde66\%，A+级）、中际旭创/新易盛（+27\textasciitilde54\%）
  \item \amark{} \textbf{等催化 / 轻仓博弈}：天准科技（7月解禁后）、柯力传感（7月Optimus投产）、五一世界/索辰科技
  \item \rmark{} \textbf{严重透支 / 规避}：绿的谐波（-84\textasciitilde-80\%）、鸣志电器（-84\%）、奥比中光（-44\textasciitilde-21\%）、中大力德/智微智能
\end{itemize}

\section{机构三维决策矩阵}

1 亿以上大资金唯一可以重仓的 5 只票：\textbf{德赛西威、协创数据、工业富联、中际旭创、新易盛}（合计建议仓位 60--70\%）。其余 13 只原则上不进入机构股票池，或单票 $\leq$5\%。

\section{未来 3--6 个月核心催化}

\begin{enumerate}[leftmargin=2em]
  \item 7 月 Optimus V3 正式投产（带动绿的谐波/柯力/奥比脉冲）
  \item 7--8 月宇树科创板 IPO（带动中大力德/鸣志）
  \item Q3 Coalition 新增 20 家国内成员（协创+仿真层弹性最大）
  \item Q4 DRIVE Thor 首批量产 10 万辆（德赛西威 240 亿营收弹性）
\end{enumerate}

\section{最大风控要点}

\begin{riskbox}[Top 3 风控要点]
\begin{enumerate}[leftmargin=1.5em]
  \item \textbf{中美禁售黑天鹅}：组合必须保留 20--30\% 与 NVIDIA 无关敞口，单一绑定标的 $\leq$15\%
  \item \textbf{7--9 月解禁潮}（90\% 概率发生）：绿的谐波 18.5\%、索辰 11.7\%、智微 6.2\%、鸣志 5.7\%、中大力德 9.3\% --- 解禁前 30 天一律减仓 50\%
  \item \textbf{追高泡沫区}：任何票高于乐观档 10\% 以上绝不追高
\end{enumerate}
\end{riskbox}
